\subsubsection{3.3 拟解决的关键科学问题}

% 破题段：必要性 → 现状不足 → 为什么是这三个 → 抛结论
海域电磁频谱侦测是海上执法与海域态势感知的重要支撑，需在海杂波强干扰、端侧算力与链路受限、多平台动态对抗等约束下，实现频谱信号的可靠感知、轻量留存与智能研判。现有方法多面向单一技术环节、相对稳定环境或资源较充足平台，各环节相对分离、缺乏全链条协同，且诸多环节仍依赖人工比对与经验研判，难以适配海杂波场景：感知端海杂波建模、传播分析与干扰对消相互分离，微弱信号难以稳定提取；存储端价值评估、表征压缩与可信存证缺乏面向连续频谱流的一体化设计，海量数据难以高效可信留存；认知端知识图谱、目标识别与协同决策多独立开展，缺乏知识动态演化与多节点可信协同，难以支撑及时可信的研判。

上述三方面分别对应海域侦测“感知--存储--认知”全链条的关键环节，前序输出构成后序输入，相互耦合、缺一不可，因此本项目不孤立处理单一环节，而将其凝练为三个层层递进的关键科学问题并置于海杂波环境下协同研究，这既是难点也是创新所在：针对微弱信号提取，凝练\textbf{科学问题一}；针对端侧频谱数据的价值评估与轻量可信存储，凝练\textbf{科学问题二}；针对多平台目标识别与可信协同决策，凝练\textbf{科学问题三}。三者的具体内涵与解决思路如下：

% 问题一
\textbf{科学问题一：海杂波干扰下微弱电磁频谱信号的传播规律表征与精准提取}

% 小帽段
在复杂海洋环境与高密度电磁对抗背景下，如何从非平稳强海杂波中高保真地提取侦测目标的微弱辐射信号，本质上是一个融合物理传播规律与小样本学习的弱信号精准提取问题。
% A
一方面，海杂波呈现非高斯、非平稳与多径时空耦合特性，其与海面环境及电磁波传播的交互机理尚未充分揭示，难以为微弱信号提取提供可解释的传播规律先验；
% B
同时，微弱信号与强杂波在时频域深度交叠，传统基于统计平稳假设的检测与杂波抑制方法在非高斯、非线性强干扰下普遍失效，残余杂波严重淹没目标信号；
% C
另一方面，侦测信号常具低信噪比、高动态与短时猝发特征且先验样本稀缺，单一域表征与纯数据驱动模型难以稳健刻画其时频结构。
% 小总结
因此，\textbf{如何揭示海杂波环境下电磁频谱信号的时空耦合传播规律，在强干扰抑制与微弱特征保真之间求取平衡，实现强海杂波背景中微弱辐射信号的精确提取，破解辐射目标“看不清”难题，是本项目亟需解决的关键科学问题}。

% 问题二
\textbf{科学问题二：有限算力端侧频谱数据的价值评估与轻量可信存储}

% 小帽段
在海上持续侦测任务中，如何在端侧有限算力下对连续产生的海量高维频谱数据流实现关键价值评估、特征保真的表征压缩与可信留存，本质上是受限资源下的频谱数据轻量存储与可信管理问题。
% A
一方面，海量高维数据与受限算力存在根本矛盾，固定阈值或均匀采样缺乏对时空关联价值与事件重要性的细粒度评估，易遗漏关键异常；
% B
同时，频谱数据连续产生、体量大、维度高，常规采样压缩易丢失关键特征，且缺乏高效索引机制，难以兼顾长效存储、特征保真与快速溯源，事后取证检索耗时久；
% C
另一方面，海上侦测环境对抗性强，传统中心化存储存在被恶意篡改的风险，难以满足执法取证对数据溯源与完整性的硬性要求，而端侧硬件又难以承载高开销的共识与加密计算。
% 小总结
因此，\textbf{如何在端侧有限算力下实现频谱数据的价值评估、特征保真的表征压缩与防篡改可追溯管理，在资源开销与证据完整可信之间求取平衡，破解海量频谱数据流留存取证“效率低”难题，是本项目需要攻克的关键科学问题}。

% 问题三
\textbf{科学问题三：基于频谱知识图谱的轻量化目标识别与可信协同决策问题}

% 小帽段
在多平台海域侦测任务中，如何在端侧受限算力与间歇通信条件下实现多源频谱的轻量化目标识别与多节点可信协同研判，本质上是受限资源与节点互信约束下的轻量化识别与分布式可信决策问题。
% A
一方面，辐射源类型多样、电磁环境动态变化，缺乏统一表征的频谱知识图谱与动态演化机制，端侧难以对多源频谱一致建模与持续积累，制约认知准确性与动态适应能力；
% B
同时，传统深度模型在边缘端受制于算力与功耗，难以兼顾多类微弱信号的高精度识别与实时响应，端侧高效轻量的识别推理机制仍待突破；
% C
另一方面，单站观测存在盲区且虚警率高，多站协同又受限于有限带宽、间歇链路与节点互不信任，集中式融合既不适应弱连接信道，也难以抵御被俘节点注入的虚假信息。
% 小总结
因此，\textbf{如何构建面向电磁态势的标准化频谱知识表征，在端侧受限算力与弱链路约束下兼顾微弱目标的识别精度、研判时效与多节点协同可信，破解大模型智能识别研判“协同难”难题，是本项目亟需解决的关键科学问题}。